% ============================================================================
% SEC03.TEX - Section 8.3: Membuat Main Menu
% ============================================================================

\section{Membuat Main Menu}

Buka scene \texttt{MainMenu} yang sudah kita buat di Bab 4.

\subsection{Desain Menu}
\begin{enumerate}
    \item Tambahkan \textbf{Canvas}.
    \item Tambahkan \textbf{Text - TextMeshPro} untuk Judul Game (\textit{Space Survivor}) yang besar di tengah atas.
    \item Tambahkan \textbf{UI > Button - TextMeshPro}. Ganti teksnya menjadi \textbf{PLAY}.
\end{enumerate}

\subsection{Script UI Manager}
Buat script \texttt{MainMenuManager.cs} pada Canvas:

\begin{lstlisting}[language={[Sharp]C}]
using UnityEngine;
using UnityEngine.SceneManagement; // Wajib untuk pindah scene

public class MainMenuManager : MonoBehaviour
{
    public void PlayGame()
    {
        // Pastikan nama scene sesuai
        SceneManager.LoadScene("Gameplay");
    }
    
    public void QuitGame()
    {
        Application.Quit();
    }
}
\end{lstlisting}

\subsection{Menghubungkan Tombol}
\begin{enumerate}
    \item Klik Button \textbf{Play}.
    \item Di komponen Button, cari bagian \textbf{On Click ()}.
    \item Klik (+).
    \item Drag objek Canvas (yang berisi script tadi) ke slot object.
    \item Pilih function: \texttt{MainMenuManager > PlayGame}.
\end{enumerate}
